\section{The Whole in the Part}

Christianity could have been easier to universalize if Christ had brought only propositions. His teaching might then be detached from his person, cleansed of the accidents of first-century Palestine, and retained as a superior wisdom. But the Gospel does not say merely that the Word delivered a message through human circumstances. It says, ``The Word became flesh and dwelt among us.''\endnote{John 1:1--18, especially v.~14. The incarnational claim governs the paragraph; it is not inferred from the preceding rhetorical contrast.}

This is why time has more than incidental importance in Christianity. Revelation occurs through words and deeds bound together: promise and fulfillment, healing and proclamation, supper and command, death and resurrection. The eternal Word does not cease to be eternal when he assumes a human nature. Neither does that nature become an appearance. Jesus works with human hands, thinks with a human mind, chooses with a human will, and loves with a human heart. John Paul II compresses the paradox: the Eternal enters time, and the Whole lies hidden in the part.\endnote{Second Vatican Council, dogmatic constitution \emph{Dei Verbum} (18 November 1965), nos.~2--4, teaches that revelation is realized through deeds and words with an inner unity and culminates in Christ, mediator and fullness; \sourceurl{https://www.vatican.va/archive/hist_councils/ii_vatican_council/documents/vat-ii_const_19651118_dei-verbum_en.html}{official English}. Second Vatican Council, pastoral constitution \emph{Gaudium et spes} (7 December 1965), no.~22, governs the confession of Christ's integral human action. John Paul II, encyclical \emph{Fides et ratio} (14 September 1998), nos.~7, 11--12, treats revelation as immersed in time and history and the Incarnation as the enduring synthesis of eternal and temporal; \sourceurl{https://www.vatican.va/content/john-paul-ii/en/encyclicals/documents/hf_jp-ii_enc_14091998_fides-et-ratio.html}{official English}. The brief ``Whole'' formulation is quoted from no.~12; official wording retains its own status.}

The doctrinal claim is exact. The things Jesus says and does are finite human realities: this sentence rather than every sentence, this touch, this journey, this death. But the personal subject who speaks, touches, travels, and suffers is the eternal Son. His humanity neither contains the divine essence nor offers one detachable fragment of a revelation completed somewhere above him. The mystery of God remains inexhaustible, while God the Son truly reveals the Father in and through his human life.

The declaration \emph{Dominus Iesus} makes this point with unusual directness. The historical words and deeds of Jesus are limited as human realities, yet they have the divine Person of the incarnate Word as their subject. Their finitude therefore does not make Christ one incomplete manifestation awaiting completion by a more universal synthesis. The one who acts in them is the Son.\endnote{Congregation for the Doctrine of the Faith, declaration \emph{Dominus Iesus} (6 August 2000), no.~6; nos.~13--15 govern the related claim of Christ's unique and universal salvific mediation; \sourceurl{https://www.vatican.va/roman_curia/congregations/cfaith/documents/rc_con_cfaith_doc_20000806_dominus-iesus_en.html}{official English}. John Paul II ratified and confirmed the declaration ``with sure knowledge and by his apostolic authority.'' Its direct use here is bounded to the historical Incarnation and universal mediation, and it is read together with \emph{Dei Verbum}, \emph{Fides et ratio}, \emph{Gaudium et spes}, and \emph{Nostra aetate}; it is not adopted as a general rhetoric for comparative religion.}

This is the Christian answer to the supposed contest between the local and the universal. Jesus is not universal because the details of his life can be discarded. His life bears universal significance because this particular human life belongs to the Person of the universal Word. ``There is one God,'' Paul writes, and ``one mediator between God and men, the man Christ Jesus, who gave himself as a ransom for all.'' The mediation is one; its scope is all.\endnote{1 Tim 2:1--6, especially vv.~4--6. The text joins God's will that all be saved to the one mediator; the article neither detaches the two claims nor infers that the divine will is fulfilled in every person irrespective of response.}

One mediation does not mean that every person has already received salvation, that every religious claim is interchangeable, or that grace reaches only those with explicit historical knowledge of Jesus. The Church holds these claims together rather than resolving one into another. She confesses Christ's unique mediation; she reveres what is true and holy in other religions; and she teaches that, because Christ died for all, the Holy Spirit offers every person, in a manner known to God, the possibility of being associated with the Paschal mystery.\endnote{\emph{Nostra aetate}, no.~2, supplies the respect owed to truth and holiness in other religions. Second Vatican Council, pastoral constitution \emph{Gaudium et spes} (7 December 1965), no.~22, joins Christ's death for all to the Spirit's universally offered possibility of association with the Paschal mystery; \sourceurl{https://www.vatican.va/archive/hist_councils/ii_vatican_council/documents/vat-ii_const_19651207_gaudium-et-spes_en.html}{official English}. \emph{Dominus Iesus}, nos.~13--15, governs the claim that this universal offer is accomplished through the one Christ. None of these sources teaches that all persons are in fact saved.}

The Cross brings the paradox to its sharpest point. A death outside Jerusalem, ``once for all,'' is proclaimed as an act ``for all.'' Its universality is not the generality of a maxim. It is the reach of a personal act. Love can be offered to many only by being actually offered by someone.\endnote{Heb 9:23--28, especially v.~26; 2 Cor 5:14--15. ``For all'' states the scope of Christ's death in the Pauline verse; it does not settle each person's reception of its saving fruit.}

If the Incarnation is true, the particular is no longer an obstacle reluctantly tolerated by universal truth. In the humanity of Christ, the particular has become the chosen place where the universal Lord gives himself.
